% ตัวอย่างเฉลยพาร์ทคณิตศาสตร์ที่ถูกต้อง 1 ข้อ — แม่แบบของ "ครบ 6 ส่วน"
% ห้าม \input ไฟล์นี้เข้าไปในเฉลยจริง เปิดอ่านอย่างเดียว
%
% รูปแบบถอดมาจาก ~/Documents/POSN.Computer/ข้อสอบเทียม/ตัวอ่างไฟล์เฉลย คณิต.pdf
% (50 ข้อ ทุกข้อมีครบ 6 ส่วน และมีอย่างน้อย 3 ขั้น)
%
% โจทย์: เลือกกรรมการ 3 คนจากสมาชิก 8 คน ได้กี่วิธี
% ตัวเลือกทั้งสี่มาจาก verify.py:
%   R[1] = C(8,3) = 56
%   D[1] = {'คูณ n กับ r ตรง ๆ': 24, 'หารด้วย r แทน r!': 112,
%           'นับลำดับด้วย (ใช้ P แทน C)': 336}
% ส่วน \whywrong ข้างล่างใช้ key ของ D[] ตรง ๆ ไม่ได้เขียนคำอธิบายใหม่

\Needspace{4\baselineskip}
\item
\ansline{ค.\ $56$}{}

\idea{โจทย์ถามจำนวนวิธีเลือก "คณะกรรมการ" ซึ่งไม่มีตำแหน่ง สลับชื่อกันแล้วได้คณะเดิม
จุดชี้ขาดของข้อนี้จึงอยู่ที่การแยกให้ออกว่า \textbf{นับลำดับหรือไม่นับลำดับ}
ถ้าไม่นับลำดับคือการจัดหมู่ ไม่ใช่การเรียงสับเปลี่ยน}

\howto
\stepa{1}{ถามตัวเองก่อนว่า "สลับตำแหน่งกันแล้วเป็นคำตอบเดิมไหม" คณะกรรมการ
          $\{A,B,C\}$ กับ $\{C,A,B\}$ คือคณะเดียวกัน จึง \textbf{ไม่นับลำดับ}
          ใช้การจัดหมู่}
\stepa{2}{เขียนสูตรรูปทั่วไปก่อนแทนค่า
          \[ \binom{n}{r}=\frac{n!}{r!\,(n-r)!} \]
          ในที่นี้ $n=8$ (คนทั้งหมด) และ $r=3$ (คนที่เลือก)}
\stepa{3}{แทนค่าแล้วตัดแฟกทอเรียลก่อนคูณ จะได้ตัวเลขเล็กพอคิดในใจ
          \[ \binom{8}{3}=\frac{8!}{3!\,5!}
             =\frac{8\times7\times6}{3\times2\times1} \]}
\stepa{4}{คิดเลขให้จบทีละขั้น ไม่ข้าม
          \[ \frac{8\times7\times6}{3\times2\times1}=\frac{336}{6}=56 \]}
\ansfinal{$56$ วิธี}

\whywrong{\quickck{$\binom{8}{3}$ ต้องเท่ากับ $P(8,3)$ หารด้วย $3!=6$ พอดีเสมอ
คำนวณ $P(8,3)=8\times7\times6=336$ ก่อน แล้วดูว่าตัวเลือกไหนคือ $336/6$
ตัดที่เหลือทิ้งได้ทันทีโดยไม่ต้องคิดใหม่}
ก. $24$ คือคูณ $n$ กับ $r$ ตรง ๆ
ข. $112$ คือหารด้วย $r$ แทนที่จะหารด้วย $r!$
ง. $336$ คือ $P(8,3)$ ของคนที่นับลำดับด้วย}

\pitfall{ใช้ $P(8,3)=336$ เพราะเผลอนับลำดับ ให้ถามตัวเองทุกครั้งว่า
\textbf{"สลับตำแหน่งกันแล้วเป็นคำตอบเดิมไหม"} ถ้าเดิม = จัดหมู่ ถ้าไม่เดิม = เรียงสับเปลี่ยน}

% สิ่งที่ตัวอย่างนี้ทำถูกและต้องทำตามทุกข้อ
%
% 1. ครบ 6 ส่วนตามลำดับ: \ansline \idea \howto \stepa \ansfinal \whywrong \pitfall
% 2. 4 ขั้น (ขั้นต่ำคือ 3) และไม่ข้ามพีชคณิต — แสดง 336/6 ตรงกลาง ไม่กระโดดไป 56
% 3. เขียนสูตรรูปทั่วไปก่อนแทนค่า ไม่ใช่โผล่ 56 มาเฉย ๆ
% 4. \quickck คือวิธีตัดตัวเลือกที่ใช้ได้จริงในห้องสอบ ไม่ใช่การคำนวณซ้ำรอบสอง
% 5. ทุกตัวลวงมีที่มาจาก D[] และคำอธิบายคือชื่อความผิดที่คำนวณไว้จริง
%    ไม่ใช่ "ข. ผิดเพราะคำนวณผิด"
% 6. เลือกวิธีที่ทำได้จริงในห้องสอบ — สูตรจัดหมู่ ไม่ใช่การไล่นับ 56 กรณี
%
% ตัวอย่างที่ผิดสำหรับข้อเดียวกัน — กระโดดขั้นแบบนี้คือสิ่งที่ต้องแก้:
%     \ansline{ค.\ $56$}{} จาก $\binom{8}{3}=56$
